\documentclass[11pt]{article}
\usepackage[margin=1in]{geometry}
\usepackage[T1]{fontenc}
\usepackage{lmodern}
\usepackage{amsmath,amssymb,amsthm,mathtools}
\usepackage{booktabs,array,longtable}
\usepackage{listings,xcolor}
\usepackage{xurl}
\usepackage[hidelinks]{hyperref}
\hypersetup{pdftitle={Rectangle Criteria for Girth Eight in Cyclic Lifted-Product Codes},pdfauthor={},pdfsubject={Lean verification, absorbing sets, expansion, and decoding}}
\newtheorem{theorem}{Theorem}[section]
\newtheorem{lemma}[theorem]{Lemma}
\newtheorem{proposition}[theorem]{Proposition}
\newtheorem{corollary}[theorem]{Corollary}
\theoremstyle{definition}
\newtheorem{definition}[theorem]{Definition}
\theoremstyle{remark}
\newtheorem{remark}[theorem]{Remark}
\newcommand{\Zr}{\mathbb Z/r\mathbb Z}
\newcommand{\Ftwo}{\mathbb F_2}
\newcommand{\HX}{H_{\mathsf X}}
\newcommand{\HZ}{H_{\mathsf Z}}
\newcommand{\GX}{G_{\mathsf X}}
\newcommand{\row}{\operatorname{row}}
\newcommand{\vol}{\operatorname{vol}}
\newcommand{\rank}{\operatorname{rank}}
\newcommand{\R}{\mathcal R}
\newcommand{\Q}{\mathcal Q}
\newcolumntype{P}[1]{>{\raggedright\arraybackslash}p{#1}}
\lstset{basicstyle=\small\ttfamily,columns=fullflexible,keepspaces=true,
  frame=single,framerule=.3pt,breaklines=true,showstringspaces=false,escapeinside={(*@}{@*)},
  literate={∀}{{$\forall$}}1 {≠}{{$\ne$}}1 {∧}{{$\land$}}1 {ℕ}{{$\mathbb N$}}1
  {→}{{$\to$}}1 {∈}{{$\in$}}1 {¬}{{$\neg$}}1 {≤}{{$\le$}}1
  {⟨}{{$\langle$}}1 {⟩}{{$\rangle$}}1}
\setlength{\emergencystretch}{2em}
\title{Rectangle Criteria for Girth Eight in Cyclic Lifted-Product Codes:\\
A Lean-Verified Answer and Decoding Consequences}
\author{}
\date{October 22, 2026}

\begin{document}
\maketitle
\begin{abstract}
Freij-Hollanti, Morris and \v Sapokait\.{e} asked for sufficient conditions,
expressed through $2\times2$ submatrices, for a cyclic lifted-product Tanner
graph to attain its maximum possible girth of eight, and for an analysis of
the resulting decoding-relevant properties. We give an exact criterion:
individual rectangle voltages must be nonzero and specified pairs of
rectangle voltages must not cancel. The characterization for the expanded
binary matrix is proved in Lean. For the degree-three case we derive a
bijection between eight-cycles and minimum $(4,4)$ absorbing sets, a stopping
set lower bound, and local syndrome bounds. A connected girth-eight family
with conductance tending to zero shows that the criterion alone does not
ensure uniform expansion. Exhaustive optimization at lift size fourteen
reduces the minimum absorbing-set count from 3,150 to 2,520 while preserving
252 physical and ten logical qubits. In an independent, paired min-sum
experiment with 100,000 errors per noise level, the selected code's recovery
failure rate falls from $4.385\%$ to $2.523\%$ at physical error probability
$4\%$. We distinguish the formal theorem, written structural proofs, exact
finite computations and statistical decoding evidence, and relate each to
its mathematical predecessors.
\end{abstract}

\section{The open question and the answer}\label{sec:question}

In Section~5 of \emph{Graphical Analysis of Lifted Product Code
Constructions}, Freij-Hollanti, Morris and \v Sapokait\.{e} explicitly ask
for \emph{``not just necessary, but sufficient conditions''} for the Tanner
graph to have \emph{``maximum possible girth of $8$''}, together with the
effects on other decoding-relevant graph properties~\cite{FHM2026}.
The question concerns the cyclic, single-base-matrix lifted-product
construction. It is presented as an open research question, rather than a
numbered conjecture. Its first part asks for a definite mathematical
criterion; its second part proposes a broader structural and experimental
program.

We answer the first part by an equivalence in the entries of the original
base matrix, not a test requiring construction of the much larger binary
parity-check matrix. The second part is addressed by precise consequences
for degree-three graphs, a limitation on expansion, and a finite optimization
and decoding study. Table~\ref{tab:scope} specifies the extent of each answer.
The degree-three conclusions concern fully populated $3\times3$ bases;
the girth theorem holds in every dimension with at least two rows and two
columns.

\begin{table}[ht]
\centering\small
\begin{tabular}{P{.24\linewidth}P{.43\linewidth}P{.23\linewidth}}
\toprule
Question & Answer & Evidence \\
\midrule
Sufficient girth conditions & A necessary and sufficient rectangle-pair criterion & General Lean proof \\
Absorbing sets & In degree three, minimum sets are exactly the four-variable eight-cycle supports & Lean bijection; written $(4,4)$ profile \\
Stopping sets and syndromes & Stopping size at least six; local cancellation bounds & Lean general lower bound and witness \\
Connectivity and expansion & Connectivity is separate; connected girth-eight graphs can have vanishing conductance & Written argument; finite checks \\
Finite design and decoding & Exact fixed-size minimum; paired decoder comparison & Exhaustive search and Monte Carlo \\
\bottomrule
\end{tabular}
\caption{The answer to the two parts of the source question, with evidence levels.}
\label{tab:scope}
\end{table}

\section{Predecessors and chronology}\label{sec:history}

Fossorier's exponent-sum criterion supplies the classical short-cycle
congruences~\cite{Fossorier2004}. Smarandache and Mitchell develop necessary
and sufficient QC girth tests through matrix products and derived
submatrices~\cite{SM2021}. Our walk-based proof specializes the classical
congruences and supplies the reduction from the lifted-product support to
the original input matrix.

Raveendran and Vasi\'c describe the product eight-cycle mechanism in
hypergraph-product graphs~\cite{RV2021}. Zhang's thesis studies cyclic
cancellation and cycle counting for lifted products under additional
base-factor girth hypotheses~\cite{Zhang2024}. Raveendran, Declercq and
Vasi\'c exhibit the eight-label cancellation directly for the same monomial
LP family~\cite{RDV2025}. Okada and Kasai use a related coordinate-support
argument and a forced eight-cycle in square-base hypergraph products with
constrained CPM lifts~\cite{OK2026}. These works provide the closest
mechanisms behind the proof below. Their stated settings and conclusions
are distinguished from our entrywise equivalence in Table~\ref{tab:history}.

The absorbing-set design perspective is established in classical coding
theory by Wang, Dolecek, Zhang and Wesel~\cite{WDZW2011} and its cycle
consistency matrix development~\cite{WDW2012}. McMillon, Beemer and Kelley
give extremal absorbing-set bounds, including the relevant girth-eight
degree-three lower bound~\cite{MBK2023}. Morris, Pllaha and Kelley develop
the quantum absorbing-set framework and its decoder interpretation
\cite{MPK2026}. Our consequences apply this framework to the characterized
LP family. Abiad and Castriota's later spectral treatment counts cycles and
absorbing configurations and bounds stopping sizes in a different
QC-GHP/GB setting~\cite{AC2026}.

\begin{table}[p]
\centering\small
\begin{tabular}{P{.20\linewidth}P{.18\linewidth}rP{.37\linewidth}}
\toprule
Work & Date used & $\Delta$ years & Contribution and relationship \\
\midrule
Fossorier~\cite{Fossorier2004} & August 2004 & $-21.9$ & Classical QC cycle congruences; ingredients for the input-matrix tests. \\
Wang et al.~\cite{WDZW2011} & May 31, 2011 & $-15.03$ & Absorbing-spectrum design and measured classical decoding gains. \\
McMillon et al.~\cite{MBK2023} & March 31, 2021 & $-5.19$ & Early online publication of extremal absorbing-set bounds. \\
Smarandache and Mitchell~\cite{SM2021} & May 7, 2021 & $-5.09$ & Classical necessary/sufficient QC girth conditions via matrix products. \\
Raveendran--Vasi\'c~\cite{RV2021} & October 14, 2021 & $-4.65$ & Product eight-cycle and quantum trapping-set analysis. \\
Zhang~\cite{Zhang2024} & March 27, 2024 & $-2.20$ & Thesis defense; LP cyclic cancellation with additional hypotheses. \\
Raveendran et al.~\cite{RDV2025} & March 10, 2025 & $-1.25$ & Same-family LP cancellation and controlled base-matrix decoding comparisons. \\
Morris et al.~\cite{MPK2026} & March 3, 2026 & $-0.27$ & Published quantum absorbing-set framework; earlier preprints precede this date. \\
Okada--Kasai, v1~\cite{OK2026} & April 30, 2026 & $-0.11$ & Related square-base support argument and forced eight-cycle. \\
Freij-Hollanti et al.~\cite{FHM2026} & June 10, 2026 & $0$ & Explicit sufficient-condition and decoding-properties question. \\
Abiad--Castriota~\cite{AC2026} & July 15, 2026 & $+0.10$ & Related spectral enumeration and stopping-bound program. \\
Okada--Kasai, v3~\cite{OK2026} & September 11, 2026 & $+0.25$ & Expanded finite-instance certificates in the square-base setting. \\
This manuscript/artifact & October 22, 2026 & $+0.37$ & Exact Lean characterization, structural consequences and finite benchmark. \\
\bottomrule
\end{tabular}
\caption{Signed chronology relative to June 10, 2026. A negative number
means preceding the question. Day-resolved entries use elapsed days divided
by $365.2425$; the August 2004 entry is month-resolved and rounded to one
decimal. The present artifact is dated 134 days after the question.
Dates identify the versions or publication milestones used, not dates of
private discovery.}
\label{tab:history}
\end{table}

The chronology separates earlier ingredients from an answer to the stated
question. In particular, the roughly $21.9$ years between Fossorier's
criterion and the question are not a claim that the full LP equivalence
was already stated in 2004. The present answer assembles those methods
into one exact input-matrix theorem with an explicit formal certificate.
Measured instead to this manuscript, the March 2025 direct LP predecessor
is 591 days ($1.62$ years) earlier, the April 2026 square-base predecessor
is 175 days ($0.48$ years) earlier, and the July 2026 spectral study is
99 days ($0.27$ years) earlier. The main answer here follows the published
question by 134 days ($0.37$ years).

\section{Construction and notation}\label{sec:model}

Let $r\ge1$ and let $B=(B_{ij})\in(\Zr)^{m\times n}$, where $m,n\ge2$.
Every entry represents the circulant permutation $P^{B_{ij}}$; exponent
zero represents an identity block. There are no missing input blocks.
Write $B^*$ for transposition followed by exponent negation. Expanding the
block matrices
\begin{equation}\label{eq:lp}
 B_{\mathsf X}=[B\otimes I_n\mid I_m\otimes B^*],\qquad
 B_{\mathsf Z}=[I_n\otimes B\mid B^*\otimes I_m]
\end{equation}
gives binary parity-check matrices $\HX,\HZ$ with
$r(m^2+n^2)$ columns and $rmn$ rows. This is the source paper's
Equation~(2.3)~\cite{FHM2026}, within the established lifted-product
framework~\cite{PK2021,RDV2025}.

Use checks $C(i,j,a)$ and variables $L(k,j,b)$, $R(i,\ell,b)$, with fibers
$a,b\in\Zr$. Adjacency in $\GX$ is exactly
\begin{align}
 L(k,j,b)&\sim C(i,j,b+B_{ik}),\label{eq:left}\\
 R(i,\ell,b)&\sim C(i,j,b-B_{\ell j}).\label{eq:right}
\end{align}
Left variables have degree $m$, right variables degree $n$, and checks
degree $m+n$. Thus $m=n=3$ gives a $(3,6)$-biregular graph. Girth is the
length of the shortest simple cycle; the explicit cycle below guarantees
that the graphs considered here are not forests.

Define the rectangle and hexagon voltages
\begin{align}
 \R(a,b;c,d)&=B_{ac}-B_{ad}+B_{bd}-B_{bc},\label{eq:rectangle}\\
 \Q(a,b,c;d,e,f)&=B_{ad}-B_{ae}+B_{be}-B_{bf}+B_{cf}-B_{cd}.
 \label{eq:hexagon}
\end{align}
All arithmetic is in $\Zr$. Direct cancellation gives
\begin{equation}\label{eq:pair}
 \Q(a,b,c;d,e,f)=\R(a,b;d,e)+\R(b,c;d,f).
\end{equation}

\section{The exact girth-eight characterization}\label{sec:main}

\begin{theorem}[Rectangle-pair characterization]\label{thm:main}
For the construction above, $g(\GX)=8$ if and only if both conditions hold:
\begin{align}
 &\R(a,b;c,d)\ne0
 &&\text{for }a\ne b,\ c\ne d;\label{eq:cond4}\\
 &\R(a,b;d,e)+\R(b,c;d,f)\ne0
 &&\text{for pairwise distinct }a,b,c\text{ and }d,e,f.
 \label{eq:cond6}
\end{align}
Equivalently, every eligible rectangle and every eligible hexagon voltage
is nonzero. No connectedness, primality or normalization assumption is
required.
\end{theorem}

The point of~\eqref{eq:cond6} is that nonzero individual rectangles do not
prevent cancellation between two rectangles. The two conditions are tests
on the original $m\times n$ matrix, including for the sparse product support
of~\eqref{eq:lp}.

\begin{lemma}[Four-cycle reduction]\label{lem:four}
The graph $\GX$ contains a four-cycle if and only if an eligible rectangle
voltage of $B$ vanishes.
\end{lemma}
\begin{proof}
A four-cycle alternates between two variables and two checks. A variable
and a neighboring projected check determine the variable fiber uniquely,
and conversely a variable fiber determines the adjacent check fiber.
Consequently distinct same-side vertices in this short cycle cannot have
the same projection. Two left variables can share checks only when their
second coordinates agree, yielding distinct row and column indices in
$B$; closure is exactly~\eqref{eq:rectangle}. Two right variables give the
transposed, negated criterion, which has the same zero set. A left and a
right variable can share only the projected check at their common
coordinate pair, and fiber uniqueness prevents them from sharing a second
check. Thus a mixed four-cycle is impossible. Conversely, lifting a zero
rectangle at any fixed unused coordinate produces four distinct vertices
and hence a genuine four-cycle.
\end{proof}

\begin{lemma}[Six-cycle reduction]\label{lem:six}
The graph $\GX$ contains a six-cycle if and only if an eligible hexagon
voltage of $B$ vanishes.
\end{lemma}
\begin{proof}
In a simple six-cycle each pair among its three variables shares a check,
and each pair among its three checks shares a variable. Fiber uniqueness
therefore forces three distinct projections on each side.

The three projected checks form a triangle whose edges are labeled by
the block type of the intervening variable. An $L$ edge preserves the
check's second coordinate; an $R$ edge preserves its first coordinate.
If both edge types occur, two edges share a type. Their coordinate
equalities, combined with the third edge, identify two projected checks,
a contradiction. Hence all three variables belong to the same block
family. In the left family, a cycle can be written
\[
 L(d,j)\sim C(a,j)\sim L(e,j)\sim C(b,j)
 \sim L(f,j)\sim C(c,j)\sim L(d,j).
\]
The fiber returns to its starting value exactly when
$\Q(a,b,c;d,e,f)=0$. The right family gives the transposed, signed
version of the same collection of tests. Conversely, distinct row and
column triples and a zero hexagon voltage lift to a simple six-cycle.
\end{proof}

The short-cycle closure calculations are the corresponding instances of
Fossorier's QC criterion~\cite{Fossorier2004}. The coordinate reduction
is related to the support argument of Okada and Kasai~\cite{OK2026};
here it yields an equivalence for arbitrary fully populated rectangular
exponent matrices.

\begin{lemma}[Uniform eight-cycle]\label{lem:eight}
For every $B$ with $m,n\ge2$, the graph $\GX$ contains a simple eight-cycle.
\end{lemma}
\begin{proof}
Choose $i_0\ne i_1$, $j_0\ne j_1$, and arbitrary $k,\ell$. The following
vertices, listed in cycle order, satisfy~\eqref{eq:left}--\eqref{eq:right}:
\[
\begin{array}{ll}
C(i_0,j_0,B_{i_0k}-B_{\ell j_0}), & L(k,j_0,-B_{\ell j_0}),\\
C(i_1,j_0,B_{i_1k}-B_{\ell j_0}), & R(i_1,\ell,B_{i_1k}),\\
C(i_1,j_1,B_{i_1k}-B_{\ell j_1}), & L(k,j_1,-B_{\ell j_1}),\\
C(i_0,j_1,B_{i_0k}-B_{\ell j_1}), & R(i_0,\ell,B_{i_0k}).
\end{array}
\]
Distinct selected rows and columns give four distinct checks, two distinct
left variables and two distinct right variables. The block families are
disjoint. Thus the cycle is simple, regardless of its fiber values.
\end{proof}

This is the known LP cancellation mechanism~\cite{RDV2025}, with explicit
fibers and a separate simplicity check. Its product antecedents include
\cite{RV2021,Zhang2024}.

\begin{proof}[Proof of Theorem~\ref{thm:main}]
Lemmas~\ref{lem:four} and~\ref{lem:six} equate the two voltage conditions
with absence of all four- and six-cycles. Bipartiteness excludes odd
cycles. Lemma~\ref{lem:eight} provides an eight-cycle, so the absence of
shorter cycles is equivalent to girth eight. Finally,
\eqref{eq:pair} rewrites the hexagon condition using only rectangle sums.
\end{proof}

\begin{remark}[A useful failed shortcut]
Over $\mathbb Z/3\mathbb Z$, the matrix
$\left(\begin{smallmatrix}0&0&0\\0&1&2\\0&2&1\end{smallmatrix}\right)$
has no zero eligible rectangle but does have a zero hexagon and a simple
six-cycle. Lean verifies this obstruction. It refutes the rectangle-only
shortcut; the source paper does not assert that shortcut.
\end{remark}

\section{Lean certificate and mathematical scope}\label{sec:lean}

The formal development uses Lean 4.36.0-rc1 and Mathlib revision
\nolinkurl{5f7b4198d22a9c82e1c698a379334fe128095c82}. The main library contains
75 audited theorem declarations. Its final matrix theorem uses Mathlib's
native simple-graph girth, together with an explicit binary expansion of
$\HX$; the voltage graph is proved equal to that Tanner graph.

The following source corollary is an exact extract from the checked
artifact, typeset using reversible TeX escapes. It expresses the sufficient-condition answer in rectangle-only
form. A companion theorem proves necessity.
\lstinputlisting{listings/source-answer.tex}

The listing's index types specialize to finite row and column sets, with
at least two elements each, and positive $r$ for the finite circulant
interpretation. Internally, the graph lemmas are stated at greater
generality. They establish projected distinctness, eliminate mixed block
patterns, construct native simple cycles, and identify the graph with the
expanded matrix. This is the formal content behind the brief final
application, not a standalone one-line proof of the underlying mathematics.

Eight supplementary declarations in \texttt{CyclicDecoding.lean} prove
the general two-error bound and concrete absorbing/stopping certificates;
three in \texttt{CyclicOptimized.lean} prove the two selected candidates'
girth. The three source-question corollaries are also supplementary.
A further 29 audited declarations in \texttt{CyclicStructure.lean} prove
the general absorbing-support/eight-cycle bijection, its counting
corollary and minimum-size exclusion, and the general stopping bound. All
observed axioms are among \texttt{propext}, \texttt{Classical.choice} and
\texttt{Quot.sound}; the development uses no admitted proofs or native
evaluation axiom. The structural results identified below are checked
for arbitrary finite check types and finite selected-variable sets.
The elementary $(4,4)$ incidence count, the companion-stabilizer
identification and the expansion argument retain written proofs.

\section{Consequences for degree-three decoding graphs}\label{sec:structure}

Fix $m=n=3$ and assume the conditions of Theorem~\ref{thm:main}.
For a variable set $S$, induced check degree means degree in the graph
on $S\cup N(S)$, including every edge from a selected variable.

\begin{definition}
An $(a,b)$ \emph{absorbing set} is a nonempty set of $a$ variables whose induced
graph has $b$ odd-degree checks, such that every selected variable has
strictly more even-degree than odd-degree check neighbors. It is
\emph{elementary} if every neighboring check has induced degree one or
two. A \emph{stopping set} is a nonempty variable set with no neighboring
check of induced degree one.
\end{definition}

These conventions agree with the source and quantum absorbing-set
framework~\cite{FHM2026,MPK2026}. The size lower bound in the next result
is consistent with the established extremal bounds of~\cite{MBK2023}.

\begin{theorem}[Minimum absorbing sets]\label{thm:absorbing}
In any variable-degree-three Tanner graph of girth eight, no absorbing set
has fewer than four variables. Its four-variable absorbing sets are
exactly the supports of simple eight-cycles, with one such cycle per
support. All are elementary $(4,4)$ sets. Consequently their multiplicity
equals the number of eight-cycles.
\end{theorem}
\begin{proof}
A selected variable needs at least two even check neighbors. For at most
three selected variables, each such check has degree two. No pair of
variables can share two checks, as that would make a four-cycle. Sets
of size one or two therefore fail absorption. A size-three set would
need three pair checks forming a six-cycle, also excluded.

Take the four variables of an eight-cycle. An additional edge from one
of its checks to another selected variable creates a four-cycle. An
outside check meeting two selected variables creates a four-cycle when
they are consecutive, or a six-cycle when they are opposite. Thus the
four cycle checks have induced degree two and the four remaining
neighbors are distinct checks of degree one. Each variable has two even
neighbors and one odd neighbor.

Conversely, an absorbing set of four variables cannot meet a degree-four
check: every other even check would share a pair with it and create a
four-cycle, leaving each variable at most one even neighbor. Nor can it
meet a degree-three check: each of that check's variables could then have
at most one even neighbor, paired with the fourth variable. Hence all
check degrees are at most two. Replace every degree-two check by an edge
between its variables. The resulting normal graph is simple, triangle-free,
has four vertices and minimum degree two. It is exactly a four-cycle.
This also proves uniqueness of the corresponding Tanner eight-cycle.
\end{proof}

Lean proves the support characterization, exclusion of smaller nonempty
absorbing sets, and a genuine equivalence between four-variable absorbing
supports and unoriented eight-cycles. A cycle is represented by its variable
and check supports with an explicit simple cyclic ordering; rotations and
reversals are not counted separately. The development constructs a native
Mathlib simple closed walk of length eight from each certificate and proves
that the variable support uniquely determines the check support. The
$4$ odd-check count in the elementary $(4,4)$ description is established
by the written incidence argument above.

The conclusion is a sharp threshold, not removal of absorbing sets:
Lemma~\ref{lem:eight} makes four-variable sets unavoidable in this model.
It converts the next design objective into eight-cycle minimization.

\begin{proposition}[Stopping size and a sharp example]\label{prop:stopping}
A Tanner graph of girth at least eight and minimum variable degree
$d\ge2$ has stopping-set size at least $2d$. For the normalized
$3\times3$ constructions here, this bound is attained at six.
\end{proposition}
\begin{proof}
Choose a variable $v$ in a stopping set and $d$ of its check neighbors.
Each check needs a second selected variable. These $d$ witnesses are
distinct, since a coincidence creates a four-cycle. Choose one witness
$a$. At least $d-1$ other checks of $a$ each need another selected
variable. The new witnesses are distinct and cannot equal $v$ by the
four-cycle exclusion; none can be a different first-level witness by
the six-cycle exclusion. There are at least $1+d+(d-1)=2d$ variables.

If the first row and column of $B$ vanish, select
\begin{equation}\label{eq:sixset}
 S_6=\{L(0,j,0):0\le j<3\}\cup\{R(i,0,0):0\le i<3\}.
\end{equation}
Exactly the nine checks $C(i,j,0)$ meet this set, each twice.
It is a stopping set of size six and has zero syndrome.
\end{proof}

The general inequality is checked by
\texttt{stopping\_card\_ge\_twice\_degree}; its hypotheses explicitly
exclude four- and six-cycles, require a nonempty stopping set, and impose
the degree lower bound only on its selected variables. The corollary
\texttt{stopping\_minimum\_six} specializes it to degree three.

The same support is a row of $\HZ$, indexed by $(0,0)$ and fiber zero.
It represents a stabilizer, so this fact does not assign quantum logical
distance six. For the binary classical kernel of $\HX$, the stopping
lower bound and~\eqref{eq:sixset} do give minimum nonzero weight six.
Lean checks the concrete induced-degree profile over $\mathbb Z/7\mathbb Z$.

\begin{proposition}[Small-error syndrome bounds]\label{prop:syndrome}
In a four-cycle-free Tanner graph, errors at two distinct variables
$v,w$ produce at least $\deg(v)+\deg(w)-2$ odd checks. In degree three
and girth at least eight, one, two and three errors expose at least
three, four and five degree-one checks, respectively.
\end{proposition}
\begin{proof}
Distinct variables share at most one check. The two-error syndrome is
$N(v)\mathbin{\triangle}N(w)$, proving the first assertion; this argument
is formalized in Lean. For $S$ with $a\le3$ variables, the induced
incidence graph is a forest, since any cycle would have length at most
six. If $E=\sum_{v\in S}\deg(v)$ and there are $c\ge1$ components,
then $|N(S)|=E-a+c\ge E-a+1$. If $t$ checks have degree one,
$E\ge t+2(|N(S)|-t)$, whence $t\ge E-2a+2$. Set $E=3a$.
\end{proof}

These are graph and syndrome statements. Failure-inducing sets also depend
on decoder dynamics, as emphasized by~\cite{RV2021,MPK2026}. Section
\ref{sec:decoding} tests that distinction directly.

\section{Connectivity and a limitation on expansion}\label{sec:expansion}

Girth eight alone does not ensure connectedness. For example, doubling
every exponent of the modulus-seven matrix below and working modulo
fourteen yields two components, because every edge preserves fiber parity.
Direct BFS gives two components of 189 vertices each.

\begin{theorem}[Connected girth eight with vanishing conductance]
Let $r=7k$, $k\ge1$, and interpret the fixed integer matrix
\begin{equation}\label{eq:baseline}
 B_0=\begin{pmatrix}0&0&0\\0&1&2\\0&3&6\end{pmatrix}
\end{equation}
modulo $r$. The corresponding graphs are connected and have girth eight.
For $k\ge2$, their edge conductance satisfies
\begin{equation}\label{eq:conductance}
 \phi(\GX)\le\frac{4}{3\lfloor r/2\rfloor}\longrightarrow0.
\end{equation}
\end{theorem}
\begin{proof}
All eligible rectangle and hexagon sums for $B_0$ are nonzero modulo seven,
as checked by Lean in the modulus-seven case. If any vanished modulo
$7k$, it would vanish modulo seven. Theorem~\ref{thm:main} thus gives
girth eight for every $k$. The upper-left rectangle has voltage one, which
generates the whole cyclic group. The source's connectivity criterion
(Theorem~3.6 of~\cite{FHM2026}) therefore gives connectedness.

Let $S$ contain all vertex types with fibers
$0,\ldots,\lfloor r/2\rfloor-1$. Define
\[
 \phi(\GX)=\min_{\varnothing\ne T\ne V}
 \frac{|\partial T|}{\min(\vol(T),\vol(V\setminus T))}.
\]
There are nine degree-six checks and eighteen degree-three variables
per fiber, or 108 units of volume. For $k\ge2$, both fiber arcs have
length at least seven. Each matching of displacement $\epsilon$, where
$|\epsilon|\le6$, crosses the cut in exactly $2|\epsilon|$ edges.
The sum of the entries of $B_0$ is 12; each appears in three left and
three right edge types. Thus $|\partial S|=2\cdot6\cdot12=144$, and
the smaller volume is $108\lfloor r/2\rfloor$. This proves
\eqref{eq:conductance}.
\end{proof}

The cuts at $r=14,28,56$ have conductance ratios $0.190476$, $0.095238$
and $0.047619$, respectively, confirmed by direct enumeration. These
are upper bounds from specified cuts. They show that the girth criterion
alone cannot provide a uniform expansion guarantee. Expansion remains
a separate design objective, as in the source's hierarchy~\cite{FHM2026}
and the related spectral program~\cite{AC2026}.

\section{Exhaustive finite optimization}\label{sec:optimization}

Fix lift size $r=14$ and normalize the first row and column to zero:
\begin{equation}\label{eq:normalized}
 B(a,b,c,d)=\begin{pmatrix}0&0&0\\0&a&b\\0&c&d\end{pmatrix}.
\end{equation}
Girth eight forces each interior entry to be nonzero, leaving
$13^4=28,561$ candidates. Applying the exact rectangle/hexagon tests
and the connectivity condition $\gcd(14,a,b,c,d)=1$ leaves 10,368 bases.
In this normalized model the four interior exponents generate the
cycle-voltage subgroup, which explains the gcd test.

The search partitions these bases into 25 orbits under row/column
permutations, normalization, multiplication by units modulo fourteen,
and transposition. For the gauge change $B'_{ij}=B_{ij}+u_i+v_j$, the
graph isomorphism shifts fibers of $C(i,j)$ by $u_i-v_j$, of $L(k,j)$
by $-v_k-v_j$, and of $R(i,\ell)$ by $u_i+u_\ell$. Transposition
exchanges block families, swaps check coordinates and negates fibers.
These identities justify counting a representative of every orbit.
The number 25 refers to this explicit symmetry group, rather than an
assertion that no further graph isomorphisms exist.

Each representative is checked for connectedness by BFS. Eight-cycle
counts are computed both by check-projection enumeration with four
distinct variable labels and by an independent depth-eight traversal
of the Tanner graph. The validation script compares both algorithms
on all 25 representatives. It also checks explicit graph bijections
for each symmetry type and repeats the search partition.

\begin{proposition}[Computed fixed-size optimum]\label{prop:optimum}
In the normalized connected girth-eight search at $r=14$, the minimum
number of four-variable absorbing sets is 2,520. Exactly 1,728 bases
attain it, forming four orbits under the stated symmetries.
\end{proposition}

Proposition~\ref{prop:optimum} is established by exhaustive finite
computation; it is not a Lean optimality theorem. The four representative
interiors are
\[
 (1,2,4,11),\quad(1,3,6,10),\quad(1,3,7,4),\quad(1,4,6,2).
\]
The first and third encode eight logical qubits. The second and fourth
preserve the baseline's ten and are the primary rate-matched candidates:
\begin{equation}\label{eq:candidates}
 B_2=\begin{pmatrix}0&0&0\\0&1&3\\0&6&10\end{pmatrix},\qquad
 B_4=\begin{pmatrix}0&0&0\\0&1&4\\0&6&2\end{pmatrix}.
\end{equation}
For both, $\rank(\HX)=\rank(\HZ)=121$ and
$K=252-121-121=10$, by exact binary elimination. Lean independently
verifies their girth eight. Their count is 20\% below the baseline's
3,150 and $4.76\%$ below the previous candidate
$B(1,2,3,8)$ with 2,646. Equal length and dimension do not imply
equal quantum distance.

\section{Measured decoding performance}\label{sec:decoding}

\subsection{Protocol and independent checks}
We use syndrome normalized min-sum with a flooding schedule,
normalization $\alpha=0.8$, at most 100 iterations and message saturation
at absolute LLR 50. The prior is $\log((1-p)/p)$. The implementation
uses fixed-point messages with twenty fractional bits, rounds
normalization magnitudes to the nearest integer with ties up, assigns
positive sign to zero messages and hard bit zero to a zero posterior.
It halts at the first syndrome match, including the initial zero estimate.
There is no post-processing or per-design tuning. The syndrome min-sum
method and LP simulation framework have prior treatments in
\cite{Raveendran2022,Raveendran2023}; our fixed implementation is
specified to make this comparison reproducible.

Errors are independent Bernoulli($p$) in one Pauli sector, with perfect
syndrome measurements. The decoder receives the syndrome and prior.
The true error and $\HZ$ are used only afterward for scoring. If
$\hat e$ is the correction, recovery succeeds exactly when
\begin{equation}\label{eq:success}
 e+\hat e\in\row_{\Ftwo}(\HZ).
\end{equation}
Thus nonzero stabilizer residuals count as success; syndrome mismatch
or a nonstabilizer residual counts as failed recovery. This is the
standard distinction between exact recovery, degeneracy and logical
miscorrection~\cite{MPK2026}. The experiment is a one-sector
code-capacity comparison, not a circuit-level or depolarizing-channel
benchmark.

An independent Python implementation evaluates the direct extrinsic
message equations and matches the C++ engine on 128 small cases,
including an exact repetition-code oracle. Every primary code is checked
on all single-qubit errors, zero error, an injected stabilizer and a
constructed zero-syndrome logical operator. These checks distinguish
the two cases in~\eqref{eq:success}. All CSS row-pair orthogonality
relations and GF(2) ranks are checked independently.

An exploratory experiment used 10,000 shared errors at each of five
probabilities with seed 20261022. The confirmation uses a new seed,
20261023, and 100,000 independent errors at each of
$p=0.01,0.02,0.04$. At each probability all four rate-matched codes
receive the same error patterns. The confirmation therefore comprises
300,000 independent patterns and 1.2 million decoder runs. Both
rate-preserving structural minima were retained; selection did not use
the confirmation outcomes.

\subsection{Results}
\begin{table}[ht]
\centering\small
\begin{tabular}{lrrrr}
\toprule
Code & $(4,4)$ sets & $p=1\%$ & $p=2\%$ & $p=4\%$ \\
\midrule
$B_0=B(1,2,3,6)$ & 3,150 & 0.033\% & 0.369\% & 4.385\% \\
$B(1,2,3,8)$ & 2,646 & 0.031\% & 0.278\% & 2.750\% \\
$B_2=B(1,3,6,10)$ & 2,520 & 0.036\% & 0.251\% & 2.523\% \\
$B_4=B(1,4,6,2)$ & 2,520 & 0.035\% & 0.278\% & 2.607\% \\
\bottomrule
\end{tabular}
\caption{Recovery failure rates, 100,000 trials per cell. Every code
has 252 physical qubits, ten logical qubits, connectedness and girth eight.}
\label{tab:decoding}
\end{table}

At $p=4\%$, $B_2$ has approximately $42.46\%$ fewer failed recoveries
than $B_0$. Wilson 95\% marginal intervals are
$[4.260\%,4.514\%]$ and $[2.428\%,2.622\%]$, respectively.
Using the shared errors, the paired difference is $-1.862$ percentage
points, with approximate pointwise 95\% interval
$[-2.011,-1.713]$. At $p=2\%$, it is $-0.118$ points with interval
$[-0.166,-0.070]$. At $p=1\%$ there are only 33 baseline and 36
$B_2$ failures; the paired interval includes zero, so this experiment
does not establish improvement there.

The two minima are not clearly distinguished: at $p=4\%$ their paired
difference is $-0.084$ points, with interval $[-0.208,+0.040]$.
Intervals are pointwise and not corrected for multiple comparisons.
We infer neither an asymptotic threshold nor a low-noise error-floor
law from these measurements.

\begin{table}[ht]
\centering\small
\begin{tabular}{lrrr}
\toprule
Code & Syndrome mismatch & Matched logical failure & Stabilizer successes\textsuperscript{*} \\
\midrule
$B_0$ & 2,981 & 1,404 & 5,834 \\
$B(1,2,3,8)$ & 2,714 & 36 & 5,849 \\
$B_2$ & 2,503 & 20 & 6,005 \\
$B_4$ & 2,565 & 42 & 6,031 \\
\bottomrule
\end{tabular}
\caption{Outcome decomposition at $p=4\%$, per 100,000 trials.
\textsuperscript{*}Successful nonzero stabilizer residuals, in addition
to exact error recovery.}
\label{tab:outcomes}
\end{table}

The improvement includes a large reduction in logical miscorrections.
Mean iterations, counting failures at the limit, decrease from 9.689
for $B_0$ to 8.857 for $B_2$. Absorbing-set count is a structural
selection objective; the exponent changes also affect algebraic and
other graph properties, so the experiment does not isolate that count
as the sole causal variable.

\subsection{Exhaustive targeted errors and the earlier design framework}
We additionally decode every exact four-variable absorbing support,
with no other erroneous variables, at each of $p=0.01,0.02,0.04$.
Every support is checked directly for its $(4,4)$ profile. The baseline
fails on 378 of 3,150 supports, or $12\%$, always through a matched
nonstabilizer residual. The previous best and both minima succeed on
every support tested. Thus an absorbing support is not automatically
failure-inducing for this min-sum decoder.

All failing baseline supports are mixed. Their residuals have weight
six in 294 cases and weight eight in 84 cases. Explicit support
witnesses are checked for zero syndrome and exclusion from
$\row(\HZ)$. The weight-six witness bounds the baseline's distance
in the tested Pauli sector by six. No exact optimized-code distance
is assigned here. Under Bernoulli noise, these disjoint exact supports
contribute $378p^4(1-p)^{248}$ to the baseline's failure probability;
they are one error subclass, not an estimate of the full failure rate.

Raveendran, Declercq and Vasi\'c already demonstrate that different
base matrices can produce different decoding performance at equal
length, rate and girth eight~\cite[Sec.~V-A, Eq.~(92), Fig.~1]{RDV2025}.
Their comparison uses 600-qubit codes and BP+OSD-10. Our baseline
also satisfies their column-partition pattern $c_0+c_2=2c_1$ modulo
fourteen; the previous best and both minima avoid every corresponding
three-column or three-row relation. This links our example to their
low-weight-logical design framework without assigning an optimized
distance or claiming a new decoder. The finite optimum and the
specified min-sum benchmark are the concrete contributions of this study.

\section{Reproducibility and remaining scope}\label{sec:artifact}

The companion artifact is organized under
\texttt{formalization/cyclic-lifts} in the anonymous project repository
\url{https://github.com/alectionary/vibe-math}. The manuscript package
contains a snapshot of all Lean sources, the pinned Lake configuration,
experiment programs, JSON results, encoded paired outcomes and validation
scripts needed for this paper. Its manifest records SHA-256 hashes;
the snapshot does not depend on a moving branch or unpublished working
tree. Build instructions are included with the package.

The library and consumers verify the main theorem and concrete
certificates using the pinned Lean/Mathlib environment. The standard-library
Python search and C++17 decoder reproduce the finite experiments.
The statistical audit reconstructs counts and paired comparisons from
stored per-trial flags. No physical experiments were performed.

The precise girth-condition question is settled by
Theorem~\ref{thm:main}. Its decoding-properties continuation receives
three complementary answers: structural consequences, a counterexample
to uniform expansion from girth alone, and an optimized finite code with
a measured recovery improvement. A general classification across
arbitrary degrees, exact optimized quantum distances, optimization over
lift sizes, and comparisons across decoders and noisy measurement models
remain separate tasks. The absorbing-support bijection and stopping lower bound are now
Lean-checked general theorems. The elementary odd-check count, the
companion-stabilizer identification and conductance argument remain
available for further formalization.

\paragraph{Acknowledgments and research provenance.}
The question is credited to Ragnar Freij-Hollanti, Kirsten D. Morris and
Patricija \v Sapokait\.{e}; their structural program guided the choice of
targets. The predecessor contributions credited in Section~\ref{sec:history}
guided the proof comparison and interpretation of decoding outcomes.
OpenAI Codex was used for literature discovery, mathematical exploration,
Lean development, experiment implementation and manuscript drafting.

\begin{thebibliography}{99}
\bibitem{FHM2026}
R. Freij-Hollanti, K. D. Morris and P. \v Sapokait\.{e},
\emph{Graphical Analysis of Lifted Product Code Constructions},
arXiv:2606.11987v1, June 10, 2026.
\url{https://arxiv.org/abs/2606.11987v1}.

\bibitem{Fossorier2004}
M. P. C. Fossorier,
\emph{Quasi-Cyclic Low-Density Parity-Check Codes From Circulant Permutation Matrices},
IEEE Transactions on Information Theory \textbf{50}(8), 1788--1793, 2004.
\href{https://doi.org/10.1109/TIT.2004.831841}{doi:10.1109/TIT.2004.831841}.

\bibitem{SM2021}
R. Smarandache and D. G. M. Mitchell,
\emph{Necessary and Sufficient Girth Conditions for Tanner Graphs of Quasi-Cyclic LDPC Codes},
arXiv:2105.03462v1, May 7, 2021.
\url{https://arxiv.org/abs/2105.03462v1}.

\bibitem{RV2021}
N. Raveendran and B. Vasi\'c,
\emph{Trapping Sets of Quantum LDPC Codes}, Quantum \textbf{5}, 562, October 14, 2021.
\href{https://doi.org/10.22331/q-2021-10-14-562}{doi:10.22331/q-2021-10-14-562}.

\bibitem{Zhang2024}
B. Zhang, \emph{On the Left-Right Lifted Product Quantum LDPC Code},
Master's thesis, Duke University, defended March 27, 2024, Section 3.3, Theorem 5.
\url{https://dukespace.lib.duke.edu/server/api/core/bitstreams/8f596ec8-b362-43da-bf61-f8c26f0c5edc/content}.

\bibitem{RDV2025}
N. Raveendran, D. Declercq and B. Vasi\'c,
\emph{On the Minimum Distances of Finite-Length Lifted Product Quantum LDPC Codes},
arXiv:2503.07567v1, March 10, 2025.
\url{https://arxiv.org/abs/2503.07567v1}.

\bibitem{OK2026}
K. Okada and K. Kasai,
\emph{High-Girth Regular Quantum LDPC Codes from Square-Base Hypergraph Products via CPM Lifts},
arXiv:2604.27817; v1 April 30, 2026; v3 September 11, 2026.
\url{https://arxiv.org/abs/2604.27817v3}.

\bibitem{WDZW2011}
J. Wang, L. Dolecek, Z. Zhang and R. Wesel,
\emph{Absorbing Set Spectrum Approach for Practical Code Design},
arXiv:1106.0057v1, May 31, 2011.
\url{https://arxiv.org/abs/1106.0057v1}.

\bibitem{WDW2012}
J. Wang, L. Dolecek and R. Wesel,
\emph{The Cycle Consistency Matrix Approach to Absorbing Sets in Separable Circulant-Based LDPC Codes},
arXiv:1208.6094v2, December 5, 2012.
\url{https://arxiv.org/abs/1208.6094v2}.

\bibitem{MBK2023}
E. McMillon, A. Beemer and C. A. Kelley,
\emph{Extremal absorbing sets in low-density parity-check codes},
Advances in Mathematics of Communications \textbf{17}(2), 465--483, 2023;
published online March 31, 2021.
\href{https://doi.org/10.3934/amc.2021003}{doi:10.3934/amc.2021003}.

\bibitem{MPK2026}
K. D. Morris, T. Pllaha and C. A. Kelley,
\emph{Absorbing sets in quantum LDPC codes},
Advances in Mathematics of Communications \textbf{23}, 234--254, 2026;
published online March 3, 2026.
\href{https://doi.org/10.3934/amc.2026023}{doi:10.3934/amc.2026023}.

\bibitem{AC2026}
A. Abiad and N. Castriota,
\emph{Spectral and Additive Combinatorial Methods for Cycles and Absorbing Sets in Lifted-Product Quantum LDPC Codes},
arXiv:2607.13666v1, July 15, 2026.
\url{https://arxiv.org/abs/2607.13666v1}.

\bibitem{PK2021}
P. Panteleev and G. Kalachev,
\emph{Degenerate Quantum LDPC Codes With Good Finite Length Performance},
Quantum \textbf{5}, 585, 2021.
\href{https://doi.org/10.22331/q-2021-11-22-585}{doi:10.22331/q-2021-11-22-585}.

\bibitem{Raveendran2022}
N. Raveendran, N. Rengaswamy, F. Rozp\k{e}dek, A. Raina, L. Jiang and B. Vasi\'c,
\emph{Finite Rate QLDPC-GKP Coding Scheme that Surpasses the CSS Hamming Bound},
Quantum \textbf{6}, 767, 2022.
\href{https://doi.org/10.22331/q-2022-07-20-767}{doi:10.22331/q-2022-07-20-767}.

\bibitem{Raveendran2023}
N. Raveendran, J. Valls, A. K. Pradhan, N. Rengaswamy, F. Garcia-Herrero and B. Vasi\'c,
\emph{Soft syndrome iterative decoding of quantum LDPC codes and hardware architectures},
EPJ Quantum Technology \textbf{10}, 45, October 19, 2023.
\href{https://doi.org/10.1140/epjqt/s40507-023-00201-1}{doi:10.1140/epjqt/s40507-023-00201-1}.
\end{thebibliography}
\end{document}
